\documentclass[11pt]{article}

\usepackage[margin=1in]{geometry}
\usepackage{amsmath,amssymb,amsthm,mathtools}
\usepackage{empheq}
\usepackage{bm}
\usepackage{enumitem}
\usepackage{booktabs}
\usepackage{hyperref}
\hypersetup{colorlinks=true,linkcolor=blue,citecolor=blue,urlcolor=blue}

% ---- theorem environments ----
\theoremstyle{plain}
\newtheorem{theorem}{Theorem}
\newtheorem{proposition}{Proposition}
\newtheorem{corollary}{Corollary}
\newtheorem{lemma}{Lemma}
\theoremstyle{definition}
\newtheorem{definition}{Definition}
\newtheorem{remark}{Remark}
\newtheorem*{recall}{Recalled result (Lee--Chirikjian)}

% ---- macros ----
\newcommand{\g}{\mathfrak{g}}
\newcommand{\gstar}{\mathfrak{g}^{*}}
\newcommand{\so}{\mathfrak{so}(3)}
\newcommand{\SO}{\mathrm{SO}(3)}
\newcommand{\SE}{\mathrm{SE}(3)}
\newcommand{\ad}{\operatorname{ad}}
\newcommand{\Ad}{\operatorname{Ad}}
\newcommand{\adstar}{\operatorname{ad}^{*}}
\newcommand{\Adstar}{\operatorname{Ad}^{*}}
\newcommand{\Lie}[1]{\mathcal{L}_{#1}}
\newcommand{\cd}{\,\mathrm{d}}          % straight d
\newcommand{\dW}{\,\mathrm{d}W}
\newcommand{\dt}{\,\mathrm{d}t}
\newcommand{\Strat}{\circ\,\mathrm{d}W}
\newcommand{\Hess}{\operatorname{Hess}}
\newcommand{\grad}{\operatorname{grad}}
\newcommand{\dv}{\operatorname{div}}
\newcommand{\R}{\mathbb{R}}
\newcommand{\J}{\mathbb{J}}
\newcommand{\Om}{\Omega}
\newcommand{\Pib}{\Pi}
\newcommand{\Efr}{E}
\newcommand{\inner}[2]{\langle #1,#2\rangle}
\newcommand{\bnabla}{\bar{\nabla}}
\DeclareMathOperator{\tr}{tr}
\DeclareMathOperator{\skewop}{skew}
\newcommand{\half}{\tfrac{1}{2}}

\title{\bf Stochastic Dynamics of Mechanical Systems on Lie Groups:\\
Force Noise versus Configuration Noise}
\author{Working draft}
\date{\today}

\begin{document}
\maketitle

\begin{abstract}
We study how randomness enters the equations of motion of a mechanical
system whose configuration evolves on a Lie group $G$ (e.g.\ $\SO$ or $\SE$),
working on the phase space (configuration $\times$ velocity). Building on the
geometric construction of Brownian motion by Lee and Chirikjian
(\texttt{arXiv:2510.19991}), we compare two ways of injecting noise:
\emph{(A)} noise entering as a random \emph{force} in the velocity/momentum
equation, so that the configuration variable is a bounded--variation path with
zero quadratic variation; and \emph{(B)} noise entering \emph{directly} into the
configuration (reconstruction) equation, so that the configuration has nonzero
quadratic variation. Because the It\^o drift is \emph{not} covariant under
nonlinear changes of variable, the two cases behave differently. We isolate the
mechanism---the non-tensorial It\^o correction is proportional to the quadratic
variation of the noise in the transformed directions---and conclude that in
Case~(A) the geometric (reconstruction) drift is unchanged by the noise, whereas
in Case~(B) a curvature-induced drift appears. Section~\ref{sec:so3} carries out
the $\SO$ computation explicitly. Conventions are flagged where they affect the
result, and claims about the reference paper are tied to its numbered results.
\end{abstract}

\tableofcontents

%======================================================================
\section{Notation, preliminaries, and recalled results}
\label{sec:prelim}
%======================================================================

Throughout, $(M,\mathbf{g})$ is a smooth $n$-dimensional Riemannian manifold with
Levi--Civita connection $\nabla$, and $G$ is an $n$-dimensional matrix Lie group
with Lie algebra $\g=T_eG$. We use the bold symbol $\mathbf{g}$ for the metric
and the plain symbol $g$ for a group element, following the reference. We adopt
the Einstein summation convention.

\subsection{Lie-group data}

For $g\in G$ and $\eta\in\g$, left translation $L_g(h)=gh$ has tangent map
$(T_eL_g)\eta = g\eta$ (matrix shorthand), giving the left-invariant vector field
$X_\eta(g)=g\eta$. The Lie bracket is
$[\eta,\zeta]=[X_\eta,X_\zeta](e)$, and $\ad_\xi\eta=[\xi,\eta]$. With the natural
pairing $\inner{\cdot}{\cdot}$ between $\gstar$ and $\g$, the coadjoint operators
are defined by $\inner{\Adstar_g\alpha}{\eta}=\inner{\alpha}{\Ad_g\eta}$ and
$\inner{\adstar_\xi\alpha}{\eta}=\inner{\alpha}{\ad_\xi\eta}$.

A choice of inner product $\inner{\cdot}{\cdot}_\g$ on $\g$ with orthonormal
basis $\{e_1,\dots,e_n\}$ extends to a \emph{left-invariant} Riemannian metric
\begin{equation}
  \mathbf{g}(v,w)=\inner{g^{-1}v}{\,g^{-1}w}_\g, \qquad v,w\in T_gG,
  \label{eq:leftinv-metric}
\end{equation}
whose left-translated frame $E_i(g)=(L_g)_*e_i=g e_i$ is orthonormal:
$\mathbf{g}(E_i,E_j)=\delta_{ij}$. This frame is globally defined on $G$.

A Lie group is \emph{unimodular} iff $\tr(\ad_\eta)=0$ for all $\eta\in\g$;
compact groups (hence $\SO$, $\mathrm{SU}(n)$), abelian groups, and the
Euclidean groups $\SE$ are all unimodular, while $\mathrm{Aff}(\R)$ is not.

\subsection{It\^o and Stratonovich calculus, and the generator}

Given smooth vector fields $X,\sigma_1,\dots,\sigma_m\in\mathfrak{X}(M)$ and
independent scalar Wiener processes $W_i$, a Stratonovich SDE on $M$,
\begin{equation}
  \cd x = X(x)\dt + \sum_{i=1}^m \sigma_i(x)\Strat_i,
  \label{eq:strat-generic}
\end{equation}
is equivalent to the It\^o SDE $\cd x=\tilde X(x)\dt+\sum_i\sigma_i(x)\dW_i$ with
the modified drift
\begin{equation}
  \tilde X = X + \half\sum_{i=1}^m \nabla_{\sigma_i}\sigma_i .
  \label{eq:ito-strat}
\end{equation}
This is Theorem~1 of the reference. The correction term
$\half\sum_i\nabla_{\sigma_i}\sigma_i$ is the manifold generalization of the
familiar Euclidean It\^o--Stratonovich term; the Christoffel part of the
covariant derivative is what makes it geometric.

The (unique) infinitesimal generator of \eqref{eq:strat-generic} acting on
$f\in C^\infty(M)$ is (Theorem~2)
\begin{equation}
  Af = X[f] + \half\sum_{i=1}^m \sigma_i\big[\sigma_i[f]\big]
     = \tilde X[f] + \half\sum_{i=1}^m \Hess_f(\sigma_i,\sigma_i),
  \label{eq:generator}
\end{equation}
where $\Hess_f(X,Y)=X[Y[f]]-(\nabla_XY)[f]$. The generator is a genuine
invariant of the diffusion: two SDEs (in different charts, or in It\^o vs.\
Stratonovich form) define the same process in law iff they share a generator.
The \emph{drift vector field}, by contrast, is invariant only in the
Stratonovich form.

\subsection{Brownian motion on manifolds and Lie groups}

Brownian motion on $M$ is the diffusion with generator $\half\Delta$, where
$\Delta$ is the Laplace--Beltrami operator. Injecting unit noise along an
orthonormal frame $\{E_i\}$ and fixing the drift to reproduce $\half\Delta$ gives
the following (Theorem~4):
\begin{equation}
  \cd x = -\half\sum_{i=1}^n \nabla_{E_i}E_i\dt + \sum_{i=1}^n E_i\Strat_i,
  \qquad\Longleftrightarrow\qquad
  \cd x = \sum_{i=1}^n E_i\dW_i .
  \label{eq:bm-manifold}
\end{equation}
The Stratonovich drift $-\half\sum_i\nabla_{E_i}E_i$ is a curvature-induced term
measuring the failure of the frame to be parallel; the It\^o form is
\emph{drift-free} when written intrinsically. Specializing to the left-invariant
frame $E_i=ge_i$ on a Lie group (Theorem~5),
\begin{equation}
  g^{-1}\cd g = \half\sum_{i=1}^n \adstar_{e_i}e_i\dt + \sum_{i=1}^n e_i\Strat_i,
  \qquad\Longleftrightarrow\qquad
  g^{-1}\cd g = \sum_{i=1}^n e_i\dW_i,
  \label{eq:bm-liegroup}
\end{equation}
where $\gstar$ is identified with $\g$ via the inner product. The identity
$\nabla_{E_i}E_i=-g\,\adstar_{e_i}e_i$ (Koszul formula on a left-invariant
metric) converts \eqref{eq:bm-manifold} into \eqref{eq:bm-liegroup}.

\begin{recall}[Corollary~1]
If $G$ is unimodular, then $\sum_i\adstar_{e_i}e_i=0$, and the Stratonovich SDE
\eqref{eq:bm-liegroup} reduces to the drift-free form
$g^{-1}\cd g=\sum_i e_i\Strat_i$.
\end{recall}

Finally, when $M$ is embedded in $\R^{\bar n}$ with orthogonal projection
$P(x)\colon\R^{\bar n}\to T_xM$, injecting noise along the projected ambient
frame $P(x)e_i$ yields (Theorem~7)
\begin{equation}
  \cd x = -\half\sum_{i=1}^{\bar n}\nabla_{Pe_i}Pe_i\dt+\sum_{i=1}^{\bar n}Pe_i\Strat_i,
  \qquad\Longleftrightarrow\qquad
  \cd x = \half H\dt + \sum_{i=1}^{\bar n}Pe_i\dW_i,
  \label{eq:bm-embedded}
\end{equation}
where $H=\sum_i \mathrm{II}(E_i,E_i)$ is the mean curvature vector and
$\mathrm{II}$ the second fundamental form. The It\^o form now carries the drift
$\half H$, purely normal to $M$; being normal, it does not enter the generator,
but it is exactly what keeps the ambient sample path on $M$.

\subsection{Phase space of a mechanical system}
\label{sec:phasespace}

A mechanical system on $G$ lives on the tangent bundle $TG$ (Lagrangian picture)
or cotangent bundle $T^*G$ (Hamiltonian picture). Using \emph{body}
(left) trivialization,
\begin{equation}
  TG \;\cong\; G\times\g, \qquad (g,\dot g)\longmapsto (g,\xi),\quad \xi=g^{-1}\dot g\in\g,
\end{equation}
and dually $T^*G\cong G\times\gstar$ with body momentum $\mu\in\gstar$. The state
is the pair (configuration, velocity) $=(g,\xi)$. Equations of motion split into
a first-order \emph{reconstruction} (kinematic) equation and a
\emph{dynamic} equation:
\begin{align}
  \text{(reconstruction)}\qquad & \cd g = g\,\xi\dt = (T_eL_g)\xi\dt, \label{eq:recon}\\
  \text{(dynamics, Euler--Poincar\'e)}\qquad & \cd \mu = \big(\adstar_\xi\mu + f\big)\dt,
      \qquad \xi=\mathbb{I}^{-1}\mu, \label{eq:EP}
\end{align}
where $\mathbb{I}\colon\g\to\gstar$ is the (reduced) inertia tensor and
$f\in\gstar$ is the reduced external force/torque. Randomness will be introduced
either into \eqref{eq:EP} (Case~A) or into \eqref{eq:recon} (Case~B).

The coadjoint term in \eqref{eq:EP} is the geometric heart of the dynamics, so we
record why it appears. It is the exact dual of the constraint that admissible
variations of the body velocity $\xi=g^{-1}\dot g$ must satisfy.

\begin{lemma}[Euler--Poincar\'e]
\label{lem:EP}
Let $\ell\colon\g\to\R$ be a reduced (left-invariant) Lagrangian and set
$\mu=\delta\ell/\delta\xi\in\gstar$. Then Hamilton's principle
$\delta\int_{t_0}^{t_1}\ell(\xi)\dt=0$, taken over curves $g(t)$ with fixed
endpoints, is equivalent to
\begin{equation}
  \frac{\cd}{\cd t}\frac{\delta\ell}{\delta\xi}=\adstar_\xi\frac{\delta\ell}{\delta\xi},
  \qquad\text{i.e.}\qquad \dot\mu=\adstar_\xi\mu .
  \label{eq:EP-free}
\end{equation}
Adding a non-conservative force $f\in\gstar$ through the Lagrange--d'Alembert
principle yields $\dot\mu=\adstar_\xi\mu+f$, which is \eqref{eq:EP}; for the
kinetic Lagrangian $\ell(\xi)=\half\inner{\mathbb{I}\xi}{\xi}$ one has
$\mu=\mathbb{I}\xi$.
\end{lemma}

\begin{proof}
Write $\delta=\partial_\epsilon|_{\epsilon=0}$ for the variation and
$\dot{(\cdot)}=\partial_t$ for a family $g(t,\epsilon)$, and set
$\xi=g^{-1}\dot g$ and $\eta=g^{-1}\delta g\in\g$, with $\eta(t_0)=\eta(t_1)=0$.
Differentiating the two definitions and using that mixed partials commute,
\begin{align*}
  \delta\xi &= -g^{-1}(\delta g)g^{-1}\dot g + g^{-1}\,\partial_\epsilon\partial_t g
             = -\eta\xi + g^{-1}\,\partial_t\partial_\epsilon g,\\
  \dot\eta  &= -g^{-1}(\dot g)g^{-1}\delta g + g^{-1}\,\partial_t\partial_\epsilon g
             = -\xi\eta + g^{-1}\,\partial_t\partial_\epsilon g,
\end{align*}
so subtracting eliminates the common second-order term and leaves the
\emph{constrained-variation identity}
\begin{equation}
  \delta\xi=\dot\eta+[\xi,\eta]=\dot\eta+\ad_\xi\eta .
  \label{eq:constrained-var}
\end{equation}
The bracket is precisely the non-commutativity of $G$; on an abelian group it
drops out. Now vary the action, using \eqref{eq:constrained-var}, the definition
$\inner{\adstar_\xi\mu}{\eta}=\inner{\mu}{\ad_\xi\eta}$, and integration by parts:
\begin{align*}
  \delta\!\int_{t_0}^{t_1}\!\ell(\xi)\dt
  &=\int_{t_0}^{t_1}\!\Big\langle \tfrac{\delta\ell}{\delta\xi},\,\delta\xi\Big\rangle\dt
   =\int_{t_0}^{t_1}\!\big\langle \mu,\,\dot\eta+\ad_\xi\eta\big\rangle\dt\\
  &=\underbrace{\big[\inner{\mu}{\eta}\big]_{t_0}^{t_1}}_{=\,0}
   +\int_{t_0}^{t_1}\!\big\langle -\dot\mu+\adstar_\xi\mu,\;\eta\big\rangle\dt .
\end{align*}
The boundary term vanishes since $\eta$ vanishes at the endpoints. As $\eta(t)\in\g$
is otherwise arbitrary, stationarity forces $-\dot\mu+\adstar_\xi\mu=0$, which is
\eqref{eq:EP-free}.
\end{proof}

\begin{remark}[Frame-transport reading]
Equivalently, $\mu$ is the spatial momentum $\pi=\Adstar_{g^{-1}}\mu$ expressed in
the moving body frame. Newton's law is $\dot\pi=f_{\mathrm{spatial}}$, and
differentiating $\pi=\Adstar_{g^{-1}}\mu$ reproduces the $\adstar_\xi\mu$ term:
it is the group version of the transport rule
$\dot{(\cdot)}_{\mathrm{space}}=\dot{(\cdot)}_{\mathrm{body}}+\ad_\xi(\cdot)$. The
dual operator $\adstar$ (rather than $\ad$) appears because momenta live in
$\gstar$ and pair with velocities in $\g$. The sign is fixed by the left (body)
trivialization; a right-invariant reduction flips it to $-\adstar_\xi\mu$. On
$\SO$, \eqref{eq:EP-free} reads $\dot\Pi=\Pi\times\Om$ (Section~\ref{sec:so3}).
\end{remark}

\medskip
\noindent\emph{The diffeomorphism-sensitivity of the It\^o drift.}
The organizing fact for what follows is the transformation law of an It\^o SDE
under a change of variables $y=\phi(x)$. If $\cd x^k=b^k\dt+\sigma^k_i\dW_i$, then
It\^o's lemma gives
\begin{equation}
  \cd y^a = \Big(\underbrace{\partial_k\phi^a\,b^k}_{\text{tensorial}}
              \;+\; \underbrace{\half\,\partial^2_{k l}\phi^a\,\big(\sigma\sigma^\top\big)^{k l}}_{\text{non-tensorial correction}}\Big)\dt
              + \partial_k\phi^a\,\sigma^k_i\dW_i .
  \label{eq:ito-transform}
\end{equation}
The first drift piece transforms like a vector field; the second does not. The
crucial observation is that the correction is \textbf{proportional to the
quadratic covariation of the noise}, $(\sigma\sigma^\top)^{kl}$. Consequently, in
any set of directions where the noise coefficient vanishes---i.e.\ where the
process has zero quadratic variation---the It\^o drift transforms tensorially and
picks up no geometric correction. This single mechanism distinguishes the two
cases below.

%======================================================================
\section{The two cases}
\label{sec:cases}
%======================================================================

\subsection{Case A: noise as a random force (zero quadratic variation in $g$)}

Here the noise enters only the dynamic equation \eqref{eq:EP}, modeling a random
force/torque. On the trivialized phase space $G\times\gstar$,
\begin{empheq}[box=\fbox]{align}
  \cd g   &= g\,\xi \dt, \qquad \xi=\mathbb{I}^{-1}\mu, \label{eq:caseA-g}\\
  \cd \mu &= \big(\adstar_\xi\mu + f(g,\mu)\big)\dt + \Sigma(g,\mu)\dW, \label{eq:caseA-mu}
\end{empheq}
with $W$ an $m$-dimensional Wiener process and $\Sigma\colon G\times\gstar\to
(\gstar)^{m}$ the (possibly state-dependent) noise coefficient. This is the
Langevin / stochastic-Hamiltonian form of the dynamics.

\begin{proposition}[Structure of the drift in Case A]
\label{prop:caseA}
For the system \eqref{eq:caseA-g}--\eqref{eq:caseA-mu}:
\begin{enumerate}[label=(\roman*),leftmargin=2.2em]
  \item The configuration $g$ has zero quadratic variation; the reconstruction
  equation \eqref{eq:caseA-g} is identical in the It\^o and Stratonovich senses.
  \item The It\^o--Stratonovich discrepancy is confined to the $\mu$-fiber and
  equals $\half(\partial_\mu\Sigma)\Sigma$. Because the momentum fiber $\gstar$
  is a vector space carried with a flat (trivial) connection in body coordinates,
  this correction is purely Euclidean; it involves \emph{no} curvature of $G$. If
  $\Sigma$ is independent of $\mu$, the correction vanishes and the two
  conventions coincide everywhere.
  \item Under any configuration reparametrization $g\mapsto\psi(g)$ lifted to
  $TG$ (in particular the intrinsic $\leftrightarrow$ ambient-embedding map on
  the configuration factor), the reconstruction drift $g\xi$ transforms as a
  genuine vector field, with no extra term. Equivalently, no curvature-induced
  or ``pinning'' drift is generated by the noise.
\end{enumerate}
\end{proposition}

\begin{proof}[Sketch]
(i) Since $\cd g=g\xi\dt$ contains no $\dW$, the path $t\mapsto g_t$ is of bounded
variation, so $[g,g]_t\equiv 0$ and the $\half\,\cd\sigma\,\dW$ term in the
It\^o--Stratonovich bridge \eqref{eq:ito-strat} vanishes for the $g$-equation.
(ii) Only $\mu$ carries $\dW$. In \eqref{eq:ito-transform} the sole nonzero block
of $\sigma\sigma^\top$ lies in the $\mu\mu$ indices, so the correction is
$\half\,\partial^2_{\mu\mu}(\cdot)\,\Sigma\Sigma^\top$, computed with the flat
connection on $\gstar$; there is no Christoffel contribution from $G$. (iii) Take
$\phi(g,\mu)=(\psi(g),\mu)$ (a configuration chart change). Then
$\partial^2_{\mu\mu}\psi=0$ and $\partial_\mu\psi=0$, so by
\eqref{eq:ito-transform} the $g$-drift correction vanishes identically and $g\xi$
transforms tensorially.
\end{proof}

The content of Proposition~\ref{prop:caseA} is that \emph{introducing force noise
does not substantially change the drift field}: the reconstruction drift is
untouched, and the dynamic drift changes at most by a flat, chart-independent
term that is absent for additive ($\mu$-independent) noise. The geometry of $G$
never couples to the noise at the level of the drift.

\subsection{Case B: noise in the configuration (nonzero quadratic variation)}

Here the noise enters the reconstruction equation, so the configuration itself
diffuses. In the purely kinematic (overdamped) limit one drops the velocity fiber
and writes, using the left-invariant frame,
\begin{empheq}[box=\fbox]{align}
  g^{-1}\cd g &= \sum_{i=1}^n e_i\Strat_i && \text{(Stratonovich, unimodular $G$),}\label{eq:caseB-strat}\\
  g^{-1}\cd g &= \sum_{i=1}^n e_i\dW_i    && \text{(It\^o, intrinsic).}\label{eq:caseB-ito}
\end{empheq}
More generally the Stratonovich drift is $\half\sum_i\adstar_{e_i}e_i$ as in
\eqref{eq:bm-liegroup}. This is exactly the configuration Brownian motion of the
reference. A hybrid mechanical version keeps \eqref{eq:EP} and adds
$\sum_i (ge_i)\Strat_i$ to \eqref{eq:recon}; the analysis of the configuration
drift is identical.

\begin{proposition}[Structure of the drift in Case B]
\label{prop:caseB}
For \eqref{eq:caseB-strat}--\eqref{eq:caseB-ito}:
\begin{enumerate}[label=(\roman*),leftmargin=2.2em]
  \item The configuration $g$ has nonzero quadratic variation
  $[\,g^{-1}\cd g,\,g^{-1}\cd g\,]_t = n\,I\,\cd t$ in the frame.
  \item The intrinsic Stratonovich drift is $\half\sum_i\adstar_{e_i}e_i$, which
  \emph{vanishes} for unimodular $G$; the intrinsic It\^o drift also vanishes.
  \item Written in the ambient embedding $x=\iota(g)\in\R^{\bar n}$, the It\^o
  drift acquires the mean-curvature term $\half H$ of \eqref{eq:bm-embedded}. This
  term is chart-dependent: it is generated purely by the non-tensorial correction
  in \eqref{eq:ito-transform}, which is now nonzero because
  $\sigma\sigma^\top\neq 0$ in the configuration directions.
\end{enumerate}
\end{proposition}

Thus in Case~B the answer to ``does noise change the drift?'' is
\emph{representation-dependent but genuinely yes}: the intrinsic Stratonovich
drift is unchanged (indeed zero for unimodular groups), while the It\^o drift in
ambient coordinates acquires a curvature-induced correction. The two statements
are reconciled by \eqref{eq:ito-transform}: the nonlinear embedding
$\iota\colon G\hookrightarrow\R^{\bar n}$ has nonzero second derivative, and it is
contracted against the (now nonzero) configuration-space quadratic variation.

\subsection{Summary of the mechanism}
The two cases are separated by a single scalar quantity---the quadratic variation
of the configuration:
\begin{center}
\begin{tabular}{@{}lll@{}}
\toprule
 & \textbf{Case A (force noise)} & \textbf{Case B (config.\ noise)}\\
\midrule
QV of configuration $g$        & $0$                                   & $\neq 0$\\
It\^o $=$ Stratonovich for $g$? & yes                                  & no\\
Curvature/pinning drift?        & none                                 & appears in ambient It\^o form\\
Intrinsic Stratonovich drift    & mechanical drift only                & $\half\sum_i\adstar_{e_i}e_i$ ($=0$ if unimodular)\\
Effect of a config.\ diffeo.\   & drift is tensorial (invariant)       & It\^o drift shifts by $\half\partial^2\iota\,(\sigma\sigma^\top)$\\
\bottomrule
\end{tabular}
\end{center}

%======================================================================
\section{Worked example: the rigid body on $\SO$}
\label{sec:so3}
%======================================================================

\subsection{Deterministic setup and conventions}

Let $R\in\SO=\{R\in\R^{3\times3}:R^\top R=I,\ \det R=1\}$ with Lie algebra
$\so=\{S:S^\top=-S\}$. The hat map $\hat{\cdot}\colon\R^3\to\so$ satisfies
$\hat a\, b=a\times b$, with inverse (vee) $\vee$. We use the body angular
velocity $\Om\in\R^3$, $\ \hat\Om=R^\top\dot R$, so the reconstruction equation
is $\dot R=R\hat\Om$.

Following the reference, the inner product on $\so$ is
$\inner{\eta}{\zeta}_\g=\half\tr(\eta^\top\zeta)=(\eta^\vee)^\top(\zeta^\vee)$,
with orthonormal basis $e_i^{\g}=\hat e_i$ ($\{e_i\}$ the standard basis of
$\R^3$) and left-invariant orthonormal frame $E_i=R\hat e_i$. Because $\SO$ is
unimodular and its induced (bi-invariant) metric coincides with the left-invariant
one, the condition \eqref{eq:leftinv-metric}$\leftrightarrow$ambient holds, so
both the intrinsic and embedded constructions apply.

The rigid-body dynamics use the inertia matrix $\J=\J^\top\succ0$ and body
angular momentum $\Pib=\J\Om$. Euler's equation with body torque $\tau$ is
\begin{equation}
  \J\dot\Om = \J\Om\times\Om + \tau
  \qquad\Longleftrightarrow\qquad
  \dot\Pib=\Pib\times\Om+\tau=\adstar_\Om\Pib+\tau,
  \label{eq:euler}
\end{equation}
where under the identification $\gstar\cong\g\cong\R^3$ we have
$\adstar_\Om\Pib=\Pib\times\Om$. \emph{Convention flag:} the sign of the gyroscopic
term follows from $\J\dot\Om+\Om\times\J\Om=\tau$; a different hat/bracket sign
convention flips it but does not affect any conclusion below.

\subsection{Case A on $\SO$: rigid body driven by a random torque}

Model the torque as $\tau\dt\to \tau(R,\Om)\dt+\Sigma(R,\Om)\dW$ with $W$ a
$3$-dimensional Wiener process. On $\SO\times\R^3$,
\begin{align}
  \cd R   &= R\hat\Om\dt, \label{eq:A-so3-R}\\
  \cd \Om &= \J^{-1}\!\big(\J\Om\times\Om+\tau\big)\dt + \J^{-1}\Sigma\dW. \label{eq:A-so3-Om}
\end{align}

\paragraph{No configuration drift from the noise.}
Equation \eqref{eq:A-so3-R} contains no $\dW$, hence $[R,R]_t\equiv0$. The
constraint is preserved without any correction:
\begin{equation}
  \cd(R^\top R)=(\cd R)^\top R+R^\top(\cd R)+\underbrace{(\cd R)^\top(\cd R)}_{=0}
  =\big(\hat\Om^\top+\hat\Om\big)R^\top R\dt = 0,
\end{equation}
using $\hat\Om^\top=-\hat\Om$. Compare this with Case~B below, where the vanishing
of $\cd(R^\top R)$ requires a nonzero It\^o drift. In particular, the ambient
matrix form of \eqref{eq:A-so3-R} has drift $R\hat\Om$ and \emph{no}
mean-curvature (pinning) term.

\paragraph{Fiber correction only.}
The only It\^o--Stratonovich discrepancy is in \eqref{eq:A-so3-Om}. Writing
$\sigma(R,\Om)=\J^{-1}\Sigma$, the Stratonovich version of \eqref{eq:A-so3-Om} has
the same drift plus $\half\sum_a(\partial_\Om\sigma_a)\sigma_a$, a Euclidean term
on the flat fiber $\R^3$. If $\Sigma$ is constant, \eqref{eq:A-so3-Om} is
simultaneously an It\^o and a Stratonovich equation. \emph{Conclusion:} the
phase-space drift is the deterministic Euler drift
$\big(R\hat\Om,\ \J^{-1}(\J\Om\times\Om+\tau)\big)$, altered at most by a flat
fiber term; the curvature of $\SO$ does not enter.

\subsection{Case B on $\SO$: configuration Brownian motion}

Now inject noise into the configuration, i.e.\ the Brownian motion of the
reference. Since $\SO$ is unimodular, the intrinsic Stratonovich form is
drift-free:
\begin{equation}
  \cd R=\sum_{i=1}^3 R\hat e_i\Strat_i. \label{eq:B-so3-strat}
\end{equation}
Convert to It\^o in the ambient $\R^{3\times3}$ coordinates using
$A\Strat=A\dW+\half\,\cd A\,\dW$ with $A_i=R\hat e_i$. Because now $R$ carries the
noise,
\begin{equation}
  \cd(R\hat e_i)\,\dW_i
  =\Big(\sum_j R\hat e_j\Strat_j\Big)\hat e_i\,\dW_i
  =\sum_j R\hat e_j\hat e_i\,\dW_j\dW_i
  =R\hat e_i^{\,2}\dt,
\end{equation}
using $\dW_j\dW_i=\delta_{ij}\dt$. Summing and using
$\sum_{i=1}^3\hat e_i^{\,2}=-2I$,
\begin{equation}
  \half\sum_{i=1}^3 \cd(R\hat e_i)\,\dW_i=\half R\Big(\sum_i\hat e_i^{\,2}\Big)\dt=-R\dt,
\end{equation}
so the It\^o form is
\begin{equation}
  \cd R=-R\dt+\sum_{i=1}^3 R\hat e_i\dW_i. \label{eq:B-so3-ito}
\end{equation}
The drift $-R$ is exactly $\half H$ with mean curvature vector $H=-2R$
(the reference's Corollary~2), and it is precisely what enforces
$\cd(R^\top R)=0$ in the It\^o calculus:
\begin{equation}
  \cd(R^\top R)=\big(\!-\!R^\top R-R^\top R\big)\dt
  +\sum_i\hat e_i^\top R^\top R\hat e_i\dt
  =\big(-2I - \textstyle\sum_i\hat e_i^{\,2}\big)\dt
  =(-2I+2I)\dt=0.
\end{equation}

\paragraph{The drift is an artifact of the embedding.}
In the \emph{intrinsic} body form the It\^o equation is drift-free,
$R^{-1}\cd R=\sum_i \hat e_i\dW_i$. The drift $-R$ appears \emph{only} after the
nonlinear embedding $\iota\colon\SO\hookrightarrow\R^{3\times3}$, whose second
derivative is contracted against the configuration quadratic variation---the
$-R$ is the concrete value of the $\half\partial^2\iota\,(\sigma\sigma^\top)$ term
in \eqref{eq:ito-transform}. This is the sense in which ``It\^o changes under
diffeomorphisms'': the same law is written with drift $0$ intrinsically and drift
$-R$ in ambient coordinates.

\subsection{Comparison and answer to the guiding question}

\begin{itemize}[leftmargin=1.4em]
  \item \textbf{Case A (random torque).} The configuration has zero quadratic
  variation, the reconstruction drift $R\hat\Om$ is unchanged by the noise, and
  no curvature/pinning term appears in \emph{any} representation. The noise
  perturbs only the velocity fiber, and only by a flat term (zero for additive
  noise). \emph{Introducing noise does not substantially change the drift field.}

  \item \textbf{Case B (configuration noise).} The configuration has nonzero
  quadratic variation. Intrinsically (body frame) the drift still vanishes---for
  the specific reason that $\SO$ is unimodular---but in the ambient embedding a
  genuine curvature-induced drift $-R=\half H$ appears. \emph{Introducing noise
  does substantially change the It\^o drift, in a chart-dependent way.}
\end{itemize}

The unifying statement is that the extra drift is proportional to the quadratic
variation of the noise in the transformed directions,
$\propto \partial^2\phi\,(\sigma\sigma^\top)$; force noise puts this variation in
the (flat) velocity fiber, while configuration noise puts it in the (curved) base,
which is why only the latter produces a geometric drift.

%======================================================================
\section*{Open points / conventions to revisit}
\addcontentsline{toc}{section}{Open points / conventions to revisit}
%======================================================================
\begin{itemize}[leftmargin=1.4em]
  \item The Euler-equation sign and the $\half$ scaling in the $\so$ inner product
  are those of the reference; a spatial-frame (right-invariant) noise model would
  replace $\adstar$ by its right-invariant analogue and is not treated here.
  \item Case~A is stated for body-frame (left-trivialized) additive/multiplicative
  torque noise. A fully intrinsic treatment of \emph{multiplicative} fiber noise
  on $T^*G$ (e.g.\ $\Sigma$ depending on $\mu$ through a non-flat fiber metric)
  would need the vertical connection made explicit; only the flat case is claimed
  above.
  \item References beyond the attached paper (kinetic Brownian motion / stochastic
  Hamiltonian systems on Lie groups) are alluded to conceptually but not cited;
  these should be pinned to specific sources before circulation.
\end{itemize}

\end{document}